\section{Complete numerical parameters}\label{app:params}
Every constant below is the literal argument of the drivers
(\texttt{examples/\allowbreak dolfinx/\allowbreak dolfinx\_\allowbreak lens.py}, \texttt{dolfinx\_\allowbreak cloak\_\allowbreak conforming\_\allowbreak vec.py},
\texttt{dolfinx\_\allowbreak panel.py}) and the analysis scripts in \texttt{examples/\allowbreak analysis/}.
\begin{table}[!tbp]
\centering\small
\caption{Complete parameter set.}\label{tab:params}
\setlength{\tabcolsep}{4pt}
\begin{tabular}{@{}lll@{}}
\toprule
quantity & lens & cloak \\ \midrule
lamina $(E_{f1},E_{f2},G_{12},\nu_{12})$, $\rho$ & \multicolumn{2}{c}{$(131,\,9,\,5,\,0.27)$, $1.6$}\\
frequency $\omega$ (design) & $40$ & $50$\\
structural damping $\eta$ & $0.02$ & $0.02$\\
domain & $4\times3$ & $8\times5$, void $R_V=0.55$ at $(3.2,2.5)$\\
mesh & structured, $130\times98$ ($\times2$ triangles) & \texttt{gmsh} conforming, $h=0.045$ ($h/2$ on rim)\\
sponge $(c_0,m)$, Eq.~\eqref{eq:sponge} & $(80,\,0.6)$ & $(70,\,0.6)$\\
source $(A,\beta,x_0)$, Eq.~\eqref{eq:source} & \multicolumn{2}{c}{$(30,\,80,\,0.1)$}\\
CS-RBF $(h_s,r)$ & $(0.33,\,0.7)$, mirror-symmetric & $(0.28,\,0.6)$, shell $R_{\rm des}=1.5$ (tapered)\\
objective region & box $(0.78L_x,L_y/2)\pm0.18$ & interior, $r>R_C=1.7$\\
MMA move & $0.10$, 70 iterations & $0.15$ adaptive, 200--240 iterations\\
curvature limit $\zeta_{\rm all}$ & none / front (Table~\ref{tab:front}) & none / $1/R_V=1.82$\\
toolpath pitch $d$ & \multicolumn{2}{c}{$0.16$}\\
\bottomrule
\end{tabular}
\end{table}

\section{Algorithm}\label{app:alg}
\begin{algorithm}[H]
\caption{Orientation design for wave control (\texttt{dolfinx})}
\label{alg:wave}
\begin{algorithmic}[1]
\Require mesh, material $\mathbf Q_f$, frequency $\omega$, damping $\eta$, sponge
  $c(\vx)$, drive $\vb$, support sites $\{\mathbf s_i\}$, radius $r$, curvature
  limit $\zeta_{\rm all}$, void set $\mathcal V$
\State \textbf{mesh:} cut each void in $\mathcal V$ from the geometry
  (\texttt{gmsh}/OCC), refining the rim, so $\mathcal T_h$ conforms to it and the
  hole boundary is traction-free \Comment{no soft void, no staircase}
\State build $\mathbf B,\mathbf B_x,\mathbf B_y$ from \eqref{eq:rbf},
  \eqref{eq:rbfderiv}; build $\mathbf S$ from $\mathbf S=\tfrac12(\mathbf I-\bm\Pi)$
\State $\hat\theta\gets\mathbf 0$ \Comment{straight-fiber baseline}
\State solve $\vU_{\rm ref}$ on a \emph{separate void-free} mesh and transfer it by
  non-matching interpolation (cloak only); record the baseline metric
\For{$\mu\in\{\mu_0,\dots,\mu_{\rm max}\}$} \Comment{penalty continuation}
  \For{$t=1,\dots,T$}
    \State $\hat\theta\gets\mathbf S\hat\theta$;\quad
           $\theta\gets\mathbf B\hat\theta$
           \Comment{symmetrize, then interpolate}
    \State assemble $\mathcal A(\vth)$ as the real $4\times4$ block
           \eqref{eq:block} \Comment{voids already absent from $\mathcal T_h$}
    \State \textbf{forward:} factor $\mathcal A=\mathbf{LU}$, solve
           $\mathcal A\vU=\vb$
    \State evaluate $J$, and $E$ (focus energy) or the far-field mismatch
    \State \textbf{adjoint:} solve $\mathcal A\T\bm\Lambda=-\partial J/\partial\vU$
           (sparse LU of $\mathcal A\T$)
    \State $d_\ell\gets\partial R/\partial\theta_\ell$ by UFL differentiation
           \eqref{eq:sens}
    \State $\bm\zeta\gets\cos\theta\,(\mathbf B_x\hat\theta)
           +\sin\theta\,(\mathbf B_y\hat\theta)$;\quad
           $\mathcal C\gets\langle\max(0,\zeta^2-\zeta_{\rm all}^2)\rangle$
    \State $\mathbf g\gets$ \eqref{eq:sens} with the curvature penalty of \S\ref{sec:curl}
           \Comment{$\mathbf S\T\mathbf B\T$ plus curvature term}
    \State $\hat\theta\gets\textsc{MMA}\bigl(\hat\theta,\mathbf g;\
           -\tfrac\pi2\le\hat\theta\le\tfrac\pi2\bigr)$
  \EndFor
  \If{$\max_c|\zeta_c|\le\zeta_{\rm all}$} \textbf{break} \Comment{feasible}
  \EndIf
\EndFor
\State \textbf{delineate toolpaths:} sign-consistent branch of
  $(\cos2\theta,\sin2\theta)$; solve $\nabla\Phi=\mathbf g$ by \texttt{lsqr}
  \eqref{eq:phase}; tows $\chi=\tfrac12+\tfrac12\cos(2\pi\Phi/d)$, masked off
  $\mathcal V$
\State \Return $\hat\theta$, $\theta$, $\vU$, gain/reduction, $\chi$
\end{algorithmic}
\end{algorithm}
\section{Detailed derivations (modeling)}
\label{app:derivmodel}
This appendix carries every modeling step in full, so each equation of
\S\ref{sec:model}--\S\ref{sec:sponge} can be re-derived and independently checked.

\subsection{Orthotropic plane-stress stiffness $\mathbf Q_f$}
In the fiber (material) frame $(1,2)$, under plane stress
($\sigma_{33}=\sigma_{13}=\sigma_{23}=0$), the orthotropic compliance relating the
engineering strain $\veps=[\varepsilon_{11},\varepsilon_{22},\gamma_{12}]\T$ to the
stress $\bm\sigma=[\sigma_{11},\sigma_{22},\sigma_{12}]\T$ is
\begin{equation}
  \veps=\mathbf S\,\bm\sigma,\qquad
  \mathbf S=\begin{bmatrix} 1/E_{f1} & -\nu_{21}/E_{f2} & 0\\[1pt]
  -\nu_{12}/E_{f1} & 1/E_{f2} & 0\\[1pt] 0 & 0 & 1/G_{12}\end{bmatrix},
\end{equation}
with the reciprocity $\nu_{12}/E_{f1}=\nu_{21}/E_{f2}$, i.e.\
$\nu_{21}=\nu_{12}E_{f2}/E_{f1}$, which makes $\mathbf S$ symmetric. The stiffness
is $\mathbf Q_f=\mathbf S^{-1}$. The upper $2\times2$ block inverts with determinant
$\Delta=\tfrac{1}{E_{f1}E_{f2}}(1-\nu_{12}\nu_{21})$, giving
\begin{equation}
  Q_{11}=\frac{E_{f1}}{1-\nu_{12}\nu_{21}},\quad
  Q_{22}=\frac{E_{f2}}{1-\nu_{12}\nu_{21}},\quad
  Q_{12}=\frac{\nu_{12}E_{f2}}{1-\nu_{12}\nu_{21}},\quad Q_{66}=G_{12},
\end{equation}
the entries of \eqref{eq:Df} (\texttt{Qf\_matrix}). The shear row decouples, hence
$Q_{16}=Q_{26}=0$.

\subsection{Rotation to the global frame and $\partial\mathbf D_f/\partial\theta$}
Let $\theta$ be the fiber angle and $(c,s)=(\cos\theta,\sin\theta)$. Rotating the
material-frame engineering strain to the global frame is $\veps_{\rm mat}=\mathbf
T(\theta)\T\veps$ with the engineering-strain rotation
\begin{equation}
  \mathbf T(\theta)=\begin{bmatrix} c^2 & s^2 & -2sc\\ s^2 & c^2 & 2sc\\
  sc & -sc & c^2-s^2\end{bmatrix},
\end{equation}
so that the global stiffness is a congruence transform of the material one,
\begin{equation}
  \mathbf D_f(\theta)=\mathbf T(\theta)\,\mathbf Q_f\,\mathbf T(\theta)\T .
  \label{eq:Dfcong}
\end{equation}
Differentiating entry-by-entry ($c'=-s$, $s'=c$),
\begin{equation}
  \mathbf T'(\theta)=\frac{\rmd\mathbf T}{\rmd\theta}
  =\begin{bmatrix} -2sc & 2sc & -2(c^2-s^2)\\ 2sc & -2sc & 2(c^2-s^2)\\
  c^2-s^2 & -(c^2-s^2) & -4sc\end{bmatrix},
\end{equation}
and by the product rule on \eqref{eq:Dfcong},
\begin{equation}
  \frac{\partial\mathbf D_f}{\partial\theta}
  =\mathbf T'\mathbf Q_f\mathbf T\T+\mathbf T\mathbf Q_f\mathbf T'^{\,\T},
  \label{eq:dDf}
\end{equation}
which is the object every orientation sensitivity is built from
(cf.\ \eqref{eq:sens}).

\subsection{Exact five-term Fourier decomposition of $\mathbf D_f(\theta)$}
\label{app:fourier}
The \texttt{dolfinx} implementation behind every result in this paper evaluates
\eqref{eq:Dfcong}--\eqref{eq:dDf} directly, as a congruence product differentiated
symbolically by UFL. The rotation nevertheless has a sharper structure, which we derive
here because it exposes the material invariants through which every orientation effect
acts (Sec.~\ref{sec:theory}) and because it gives an equivalent closed form, cheaper to
evaluate, that a structured-grid implementation can use: \emph{$\mathbf D_f(\theta)$ is a
trigonometric polynomial of degree four in $\theta$, exactly}.

Start from the entries of $\mathbf T(\theta)$. Each is a quadratic monomial in
$(c,s)$ --- one of $c^2$, $s^2$, or $sc$ --- and the double-angle identities write
every such monomial in the fixed basis $\{1,\cos2\theta,\sin2\theta\}$:
\begin{equation}
  c^2=\tfrac12(1+\cos2\theta),\qquad
  s^2=\tfrac12(1-\cos2\theta),\qquad
  sc=\tfrac12\sin2\theta .
  \label{eq:double}
\end{equation}
So every entry of $\mathbf T$ --- and hence of $\mathbf T\T$ --- lives in
$\operatorname{span}\{1,\cos2\theta,\sin2\theta\}$. The congruence
\eqref{eq:Dfcong} multiplies two $\mathbf T$-factors (with the constant $\mathbf
Q_f$ between them), so each entry of $\mathbf D_f$ is a sum of \emph{products of two}
such first-harmonic terms. The product-to-sum identities collapse those products to
the second harmonic and no higher,
\begin{equation}
  \cos2\theta\cos2\theta=\tfrac12(1+\cos4\theta),\quad
  \sin2\theta\sin2\theta=\tfrac12(1-\cos4\theta),\quad
  \cos2\theta\sin2\theta=\tfrac12\sin4\theta ,
\end{equation}
so nothing beyond $4\theta$ can appear. Every entry of $\mathbf D_f(\theta)$ therefore
lies in the five-dimensional space
$\operatorname{span}\{1,\cos2\theta,\sin2\theta,\cos4\theta,\sin4\theta\}$, and
collecting the coefficient of each basis function into its own constant $3\times3$
matrix gives the \emph{exact} decomposition
\begin{equation}
  \boxed{\;\mathbf D_f(\theta)=\sum_{s=0}^{4}c_s(\theta)\,\mathbf C_s\;},\qquad
  \mathbf c(\theta)=\bigl[\,1,\ \cos2\theta,\ \sin2\theta,\ \cos4\theta,\ \sin4\theta\,\bigr],
  \label{eq:dffourier}
\end{equation}
in which the five $\mathbf C_s$ depend only on $\mathbf Q_f$ (each symmetric, since
$\mathbf D_f$ is). They are recovered once, offline, by projecting $\mathbf
D_f(\theta)$ onto the basis $\mathbf c$ over a dense $\theta$-sample --- an exact fit,
not an approximation, because the five functions span the range; in the
implementation the residual of the least-squares recovery is
$8.5\times10^{-16}$, i.e.\ machine zero, which is the numerical confirmation that the
degree is really four and not higher.

Two consequences make this the preferred form. First, the derivative no longer needs
the product rule \eqref{eq:dDf}; it is just the (elementwise) derivative of the
scalar basis,
\begin{equation}
  \frac{\partial\mathbf D_f}{\partial\theta}=\sum_{s=0}^{4}c_s'(\theta)\,\mathbf C_s,
  \qquad
  \mathbf c'(\theta)=\bigl[\,0,\ -2\sin2\theta,\ 2\cos2\theta,\ -4\sin4\theta,\ 4\cos4\theta\,\bigr],
  \label{eq:dffourierderiv}
\end{equation}
which is exact to round-off and shares the constant $\mathbf C_s$ with the value
\eqref{eq:dffourier}, so value and gradient cost one small matrix combination each.
Second, the decomposition passes \emph{through} the element integral. The fiber
element stiffness is $\mathbf k_f(\theta)=\int_e \mathbf B\T\mathbf D_f(\theta)\,
\mathbf B$, and because the $c_s(\theta)$ are scalars they factor out of the integral,
\begin{equation}
  \mathbf k_f(\theta)=\sum_{s=0}^{4}c_s(\theta)\,\underbrace{\int_e \mathbf B\T
    \mathbf C_s\,\mathbf B}_{\textstyle \mathbf k_f^{(s)}}
  =\sum_{s=0}^{4}c_s(\theta)\,\mathbf k_f^{(s)} ,
  \label{eq:kffourier}
\end{equation}
so an implementation can cache five fixed matrices $\mathbf k_f^{(s)}$ per element (for
a linear triangle, $\mathbf k_f^{(s)}=A\,\mathbf B\T\mathbf C_s\mathbf B$, since
$\mathbf B$ is constant), and every value or $\theta$-derivative of the element stiffness
is a five-term weighted sum of them. The two forms give identical stiffness and gradients;
the symbolic form is what we use.

\subsection{Isotropy annihilation (proof of Remark~\ref{rem:iso})}
$\mathbf Q_f$ is isotropic iff $Q_{11}=Q_{22}$ and $Q_{66}=\tfrac12(Q_{11}-Q_{12})$.
Then $\mathbf Q_f=\lambda\mathbf 1\mathbf 1\T+2\mu\,\mathbf P$ with $\mathbf
1=[1,1,0]\T$ and $\mathbf P$ the deviatoric projector, both invariant under the
in-plane rotation $\mathbf T(\theta)$ (an isotropic tensor is unchanged by
rotation): $\mathbf T\mathbf Q_f\mathbf T\T=\mathbf Q_f$ for all $\theta$. Hence
$\mathbf D_f(\theta)\equiv\mathbf Q_f$ and, by \eqref{eq:dDf},
$\partial\mathbf D_f/\partial\theta\equiv\mathbf 0$; by \eqref{eq:sens} every
orientation sensitivity vanishes. Orientation is a dead variable exactly when the
material is isotropic. $\square$

\subsection{Elastodynamics, time-harmonic reduction, damping, sponge}
Momentum balance for the displacement $\vu(\vx,t)$ with mass density $\rho$ and
the plane-stress constitutive law $\bm\sigma=\mathbf D_f(\theta)\veps(\vu)$ is
$\nabla\!\cdot\!\bm\sigma=\rho\,\ddot{\vu}$. Seeking time-harmonic motion
$\vu(\vx,t)=\Real\{\vu(\vx)e^{+\ii\omega t}\}$ turns $\ddot{\vu}\to-\omega^2\vu$ and $\dot\vu\to\ii\omega\vu$.
Structural (hysteretic) damping multiplies the stiffness by $(1+\ii\eta)$, and a
mass-proportional absorbing sponge $c(\vx)\ge0$ adds a term $+\ii\omega c\,\vu$
(the $\dot{\vu}\to\ii\omega\vu$ of a viscous force $c\dot\vu$), giving the strong
form \eqref{eq:strong}. Multiplying by a test function $\vv$ and integrating by
parts (the sponge edges carry the outgoing radiation, the driven edge the source)
yields the weak form and, after $P_1$ interpolation, the complex system
\eqref{eq:complex} with
$\mathbf K=\int\veps(\cdot)\T\mathbf D_f\veps(\cdot)$,
$\mathbf M=\int\rho(\cdot)\!\cdot\!(\cdot)$,
$\mathbf C=\int c(\cdot)\!\cdot\!(\cdot)$ (\texttt{\_forms}).

\subsection{Exact real four-field split}
Write the complex operator of \eqref{eq:complex} as
$\mathbf D=(1+\ii\eta)\mathbf K-\omega^2\mathbf M+\ii\omega\mathbf C=A+\ii B$ with
\begin{equation}
  A=\mathbf K-\omega^2\mathbf M,\qquad B=\eta\mathbf K+\omega\mathbf C,
\end{equation}
both real symmetric. Splitting $\vu=\vu_r+\ii\vu_i$ and $\mathbf f=\mathbf
f_r+\ii\mathbf f_i$, the equation $(A+\ii B)(\vu_r+\ii\vu_i)=\mathbf f_r+\ii\mathbf
f_i$ separates into real and imaginary parts,
\begin{equation}
  A\vu_r-B\vu_i=\mathbf f_r,\qquad B\vu_r+A\vu_i=\mathbf f_i,
\end{equation}
i.e.\ the real block system \eqref{eq:block} on
$\vU=(u_{xr},u_{yr},u_{xi},u_{yi})$. It is exactly $[\Real,-\Imag;\Imag,\Real]$ of
$\mathbf D$, so a real \texttt{PETSc} build solves the complex problem without
complex scalars.

\section{Detailed derivations (optimization)}
\label{app:derivopt}

\subsection{Discrete adjoint and design sensitivity}
Let the assembled forward problem be $\mathcal A(\vth)\vU=\vb$ with
$\mathcal A=[A,-B;B,A]$, and let $J(\vU)$ be an objective. Form the Lagrangian with
a multiplier $\bm\Lambda$,
$\mathcal L=J(\vU)+\bm\Lambda\T(\mathcal A\vU-\vb)$; on the constraint manifold
$\mathcal L=J$. Its total derivative in a design $\theta_\ell$ is
\begin{equation}
  \frac{\rmd J}{\rmd\theta_\ell}
  =\frac{\partial J}{\partial\vU}\T\frac{\rmd\vU}{\rmd\theta_\ell}
  +\bm\Lambda\T\Bigl(\mathcal A\frac{\rmd\vU}{\rmd\theta_\ell}
  +\frac{\partial\mathcal A}{\partial\theta_\ell}\vU\Bigr)
  =\Bigl(\frac{\partial J}{\partial\vU}+\mathcal A\T\bm\Lambda\Bigr)\T
  \frac{\rmd\vU}{\rmd\theta_\ell}
  +\bm\Lambda\T\frac{\partial\mathcal A}{\partial\theta_\ell}\vU .
\end{equation}
Choosing $\bm\Lambda$ to annihilate the bracket gives the adjoint equation and
reduces the sensitivity to the explicit design dependence of the residual
$R(\vth)=\bm\Lambda\T\big(\mathcal A(\vth)\vU-\vb\big)$ (with $\vU,\bm\Lambda$ frozen):
\begin{equation}
  \mathcal A\T\bm\Lambda=-\frac{\partial J}{\partial\vU},\qquad
  \frac{\rmd J}{\rmd\theta_\ell}=\bm\Lambda\T
  \frac{\partial\mathcal A}{\partial\theta_\ell}\vU .
\end{equation}
Only $\mathbf K(\vth)$ (hence $A,B$) depends on $\theta$, and only through
$\partial\mathbf D_f/\partial\theta$ \eqref{eq:dDf}; expanding $\bm\Lambda=(\bm
\lambda_r,\bm\lambda_i)$ block-wise and using $\partial A/\partial\theta_\ell=
\partial\mathbf K/\partial\theta_\ell$, $\partial B/\partial\theta_\ell=\eta\,
\partial\mathbf K/\partial\theta_\ell$ gives the explicit cell integral
\eqref{eq:sens}. The mass, sponge and source terms carry no $\theta$ and drop out,
so the gradient is entirely built from $\partial\mathbf D_f/\partial\theta$ ---
which vanishes for an isotropic material, making Remark~\ref{rem:iso} a statement
about the gradient itself. For a quadratic objective $J=\int w|\vu|^2$ the adjoint
load is linear, $\partial J/\partial\vU=2\,\mathrm{diag}(w)\vU$ (a signed weight $w$ gives the
multi-region objectives of Sec.~\ref{sec:panel}); in the code both the
adjoint load and the residual derivative are single \texttt{ufl.derivative} calls
(\texttt{\_sensitivity}), FD-verified to $\sim\!10^{-6}$.

\subsection{CS-RBF chain rule and its analytic derivatives}
With $\theta=B\hat\theta$ \eqref{eq:rbf}, the reduced gradient is
$\partial J/\partial\hat\theta=B\T(\partial J/\partial\theta)$. For the curvature
constraint we also need $\partial_x\theta,\partial_y\theta$. Writing
$W_{ci}=\phi(q_{ci})$ and $S_c=\sum_jW_{cj}$ so $B_{ci}=W_{ci}/S_c$, the chain rule
on $q_{ci}=\lVert\vx_c-\mathbf s_i\rVert/r$ with $\phi'(q)=-20q(1-q)^3$ gives
$\partial_x W_{ci}=\phi'(q_{ci})\,\partial_x q_{ci}
=-\tfrac{20(1-q_{ci})^3}{r^2}(x_c-s_{ix})$, and the quotient rule on $B_{ci}$
yields \eqref{eq:rbfderiv}, i.e.\ $\mathbf B_x=S^{-1}W_x-\mathrm{diag}(S^{-1}W_x
\mathbf 1)\,B$ (and $\mathbf B_y$ likewise); the second term is the derivative of
the partition-of-unity normalization, without which $\sum_i(\mathbf B_x)_{ci}\neq0$
and the curvature would be biased (\texttt{csrbf\_operators}).

\subsection{Fiber-curvature functional and its Jacobian}
The fiber-path curvature is the directional derivative of $\theta$ along the fiber
$\vn=(\cos\theta,\sin\theta)$,
$\zeta=\cos\theta\,\theta_x+\sin\theta\,\theta_y$ with $\theta_x=\mathbf
B_x\hat\theta$, $\theta_y=\mathbf B_y\hat\theta$. Differentiating w.r.t.\
$\hat\theta$ (product rule, $\partial\theta/\partial\hat\theta=B$),
\begin{equation}
  \frac{\partial\zeta}{\partial\hat\theta}
  =\mathrm{diag}(-\sin\theta\,\theta_x+\cos\theta\,\theta_y)\,B
  +\mathrm{diag}(\cos\theta)\,\mathbf B_x+\mathrm{diag}(\sin\theta)\,\mathbf B_y .
\end{equation}
The penalty $\mathcal C=\langle\max(0,\zeta^2-\zeta_{\rm all}^2)\rangle$ has gradient
$\partial\mathcal C/\partial\hat\theta=\tfrac{2}{n}(\partial\zeta/\partial\hat\theta)\T
[\zeta\odot\mathbf 1_{\{\zeta^2>\zeta_{\rm all}^2\}}]$ (restricted to solid cells
for the cloak), FD-verified to $\sim\!10^{-10}$ (\texttt{curl\_penalty}). It is
added to the objective with a continuation weight $\mu$ ramped until the
constraint is met.

\subsection{Phase-field toolpath extraction}
The tow field is the level set of $\Phi$ with $\nabla\Phi=\mathbf g=[-\sin\theta,
\cos\theta]\T$ (unit, orthogonal to the fiber). This overdetermined system
(two equations, one unknown) is solved in the least-squares sense
$\min_\Phi\int\lVert\nabla\Phi-\mathbf g\rVert^2$. Its Euler--Lagrange condition is
the Poisson equation $\nabla^2\Phi=\nabla\!\cdot\mathbf g$ with the \emph{natural}
Neumann data $\partial\Phi/\partial n=\mathbf g\!\cdot\!\vn$ on $\partial\Omega$.
The boundary flux is essential: dropping it (homogeneous Neumann) leaves only
$\nabla^2\Phi=\nabla\!\cdot\mathbf g$, whose solution for a nearly divergence-free
$\mathbf g$ collapses to $\Phi\approx\mathrm{const}$ ($|\nabla\Phi|\!\to\!0$), i.e.\
washed-out tows. We therefore assemble the discrete forward-difference gradient
operator $G=[G_x;G_y]$ on the grid and solve $\min_\Phi\lVert G\Phi-\mathbf
g\rVert^2$ directly (\texttt{lsqr}, one pinned node fixing the constant), which
enforces $|\nabla\Phi|\approx1$, and hence a constant tow pitch, where the director field is nearly divergence-free. The tows are the
peaks of $\chi=\tfrac12+\tfrac12\cos(2\pi\Phi/d)$ (\texttt{fiber\_phase}). A
seam-free director branch is obtained beforehand by region-growing sign
propagation (\texttt{\_orient\_consistent}).

\subsection{Symmetry projector}
The map $\bm\Pi$ exchanging each support site with its mirror image is an involutive
permutation ($\bm\Pi^2=\mathbf I$, $\bm\Pi\T=\bm\Pi$), so with
$\mathbf S=\tfrac12(\mathbf I-\bm\Pi)$,
$\mathbf S\T=\mathbf S$ and
$\mathbf S^2=\tfrac14(\mathbf I-2\bm\Pi+\bm\Pi^2)=\tfrac12(\mathbf I-\bm\Pi)
=\mathbf S$. Thus $\mathbf S$ is the \emph{orthogonal} projector onto
$\{\hat\theta:\bm\Pi\hat\theta=-\hat\theta\}$, the mirror-symmetric fiber fields.
Being self-adjoint, the same operator symmetrizes the design and the gradient, so
imposing symmetry introduces no inconsistency between the two --- the projected
gradient is exactly the gradient of the restricted problem. \qed

\subsection{MMA subproblem}
Each iteration the Method of Moving Asymptotes \cite{svanberg1987} replaces $J,\mathcal C$ by convex
separable approximations built from moving asymptotes $L_\ell<\hat\theta_\ell<U_\ell$,
\begin{equation}
  \tilde J(\hat\theta)=r+\sum_\ell\Bigl(\frac{p_\ell}{U_\ell-\hat\theta_\ell}
  +\frac{q_\ell}{\hat\theta_\ell-L_\ell}\Bigr),\quad
  \begin{cases} p_\ell=(U_\ell-\hat\theta_\ell^{(k)})^2\,(\partial_\ell J)^+\\
  q_\ell=(\hat\theta_\ell^{(k)}-L_\ell)^2\,(\partial_\ell J)^-\end{cases}
\end{equation}
($(\cdot)^\pm$ the positive/negative parts), with the asymptotes contracted or
expanded by the sign history of $\hat\theta_\ell$ and a move limit per step. The
separable convex subproblem is solved by its (cheap) dual. Because the raw
objectives are $O(10^{-4})$ we feed MMA the scale-invariant log-gradient
$g_\ell=(\partial_{\theta_\ell}J)/J$ for the lens focus, restoring $O(1)$
sensitivities; the cloak uses the raw gradient (well scaled). Bounds are
$\hat\theta_\ell\in[-\tfrac\pi2,\tfrac\pi2]$ (\texttt{cfrp\_optimizer.py : MMA}).

\section{Multifunctional panel objective}\label{app:panelobj}

With the region energies $E_R=\int_{\Omega_R}|\mathbf u|^2\,\mathrm d\mathbf x$, the four
wave jobs compose into a single scalar
\begin{equation}
  \mathcal J_{\rm wave}=W_S\log\frac{E_S}{E_S^0}
  -\lambda_P\frac{E_P}{E_P^0}-\lambda_F\frac{E_F}{E_F^0}-\lambda_O\frac{E_O}{E_O^0},
  \label{eq:panelobj}
\end{equation}
with the regions made \emph{disjoint} by priority (rims $>$ payload $>$ sensor $>$
outside) so no cell is rewarded and penalized at once, and each term divided by its
straight-fiber baseline $E_R^0$ so all four are $O(1)$ and directly comparable. Two
choices matter. The delivery reward is \emph{logarithmic}: with a plain ratio a bright
sensor term would swamp the rim penalty and the optimizer would buy delivery by letting
the fasteners get busier, whereas $\log$ has diminishing returns so protecting the rims
stays worthwhile however bright the port already is. The confinement term penalizes
interior energy \emph{outside} the target regions, which forces a clean routed beam
rather than a diffraction spray. The gradient still needs one adjoint solve per state:
for $\mathcal J_{\rm wave}=\sum_R c_R E_R$ (the log term contributing $c_S=W_S/E_S$,
re-evaluated each iterate) the adjoint load is the summed region weight
$w=\sum_R(\partial\mathcal J/\partial E_R)\,\mathbf 1_{\Omega_R}$, and the
discrete-adjoint sensitivity of \S\ref{sec:adj} returns
$\mathrm d\mathcal J_{\rm wave}/\mathrm d\theta_\ell$; the CS-RBF transpose then gives the
reduced gradient. The panel is mirror-symmetric about $y=L_y/2$, enforced exactly by a
symmetric support grid and the antisymmetric projector of \S\ref{sec:rbf}.

\begin{table}[!tbp]
\centering\footnotesize
\caption{Parameters of the multifunctional panel. Band design:
\texttt{dolfinx\_panel.py --band 44,47,50,53,56 --lam-p 3 --stiff 0.3}; the single-frequency
design uses \texttt{--band 50} with the same settings.}\label{tab:panelparams}
\begin{tabular}{@{}ll@{}}
\toprule
plate, lamina, damping & $7\times4$; lamina of Table~\ref{tab:params}; $\eta=0.02$\\
mesh & \texttt{gmsh} conforming, $h=0.055$ ($h/2$ on hole rims)\\
sponge $(c_0,m)$ & $(80,\,0.6)$ on the right, top and bottom edges\\
drive & $55\,e^{-80(x-0.15)^2}e^{-((y-2)/0.95)^2}$ ($x$-directed beam)\\
fastener holes & radius $0.20$ at $(2.55,1.05)$, $(2.55,2.95)$, $(4.35,1.05)$, $(4.35,2.95)$\\
rim regions $F$ & annuli of width $0.18$ around the holes\\
payload $P$ & square $(3.2,2.0)\pm0.5$\\
sensor $S$ & square $(5.9,2.0)\pm0.3$\\
weights $(W_S,\lambda_P,\lambda_F,\lambda_O)$ & $(1,\,3,\,1,\,0.5)$ (strong shield)\\
stiffness term & clamped at the four bolt rims, unit downward load over $P$; $w_c=0.3$\\
band & $\omega\in\{44,47,50,53,56\}$ (single-frequency design: $\omega=50$)\\
tow limit & $R=0.16$ ($\zeta_{\rm all}=6.25$); achieved $\max|\zeta|=\WPanZmax$\\
optimizer & MMA, 110 iterations from straight fibers\\
CS-RBF & mirror-symmetric support grid, as for the lens\\
\bottomrule
\end{tabular}
\end{table}

\paragraph{Operational bandwidth.}
A single-frequency optimum is a single-tone trick. We make the design robust across an
operational band $[\WPanBLo,\WPanBHi]$ by averaging the objective over $\WPanNF$ sample
frequencies,
\begin{equation}
  \mathcal J=\frac1K\sum_{k=1}^{K}\mathcal J_{\rm wave}(\omega_k;\theta),\qquad
  \frac{\mathrm d\mathcal J}{\mathrm d\theta_\ell}
  =\frac1K\sum_k\frac{\mathrm d\mathcal J_{\rm wave}(\omega_k)}{\mathrm d\theta_\ell},
  \label{eq:panelband}
\end{equation}
one forward and one adjoint solve per frequency per iterate. The band optimization earns
its keep at the band \emph{edges}: a design tuned at the center alone holds the shield and
delivery across the band but lets the fastener protection degrade toward the edges, which
the averaged design restores (Figure~\ref{fig:panel}d).

\paragraph{Stiffness.}
The same orientation is scored for static stiffness. The panel is clamped at the four bolt
rims (the bolted mount) and carries a service load over the payload footprint; the
compliance $c=\mathbf f^\top\mathbf u$ and its self-adjoint orientation sensitivity
$\mathrm d c/\mathrm d\theta_\ell=-\int_{\Omega_\ell}
\boldsymbol\varepsilon(\mathbf u)^\top(\partial\mathbf D_f/\partial\theta)
\boldsymbol\varepsilon(\mathbf u)\,\mathrm d\mathbf x$ --- one extra linear solve, reusing
the static factorization --- enter the objective as $-w_c\,c/c_0$. The fiber angle is largely
spent on the wave, so with this weight ($w_c=0.3$, an explicit term of the objective) the
band-optimized layout reaches $\WPanStiff\times$ the straight-fiber stiffness at no added
mass; pushing stiffness
harder trades directly against the routing and shielding. A density/topology field would
decouple stiffness and mass from the wave, but a SIMP soft-void contradicts the
conforming-void modeling used throughout and is not pursued here.

\section{Numerical verification and reproducibility}\label{app:verify}
\paragraph{Software, solver and files.} Table~\ref{tab:env} lists the software and solver settings. Every
linear solve is a sparse direct LU factorization (MUMPS through PETSc), so there is no iterative tolerance;
the forward and adjoint residuals are at round-off, and the discrete power balance \eqref{eq:balance} closes to
$\WPowBal$. Every quoted number is a macro generated from the result files by
\texttt{wt\_\allowbreak paper\_\allowbreak assets.py} (\texttt{docs/paper/wt\_numbers.tex}), and every design
shown is stored with its mesh centroids and orientation in \texttt{results/} (e.g.\
\texttt{dolfinx\_lens\_data.npz}, \texttt{lens\_absorber.npz}, \texttt{cloak\_headtohead.npz},
\texttt{dolfinx\_panel\_band\_repro\_data.npz}).

\begin{table}[!tbp]
\centering\footnotesize
\caption{Software environment and solver settings.}\label{tab:env}
\begin{tabular}{@{}ll@{}}
\toprule
FE library & \texttt{dolfinx} 0.7.3, \texttt{ufl} 2023.2.0, \texttt{basix} 0.7.0\\
linear algebra & PETSc 3.20.6, real scalars (float64); MUMPS sparse direct LU (forward and adjoint)\\
meshes & structured triangles (lens, $n_x=130$); \texttt{gmsh} 4.15 OCC, conforming holes, rim refinement $h/2$\\
elements & $[\mathbb P_1]^4$ state (real four-field split), $\mathbb{DG}_0$ orientation\\
optimizer & MMA \cite{svanberg1987}, move limit 0.10 (lens), 0.15 (cloak); iterations in Table~\ref{tab:params}\\
other & NumPy 1.26.4, SciPy 1.15.3; figures with Matplotlib (pgf backend)\\
\bottomrule
\end{tabular}
\end{table}

\paragraph{Checks.} The checks are produced by the scripts \texttt{wt\_\allowbreak verify.py},
\texttt{cloak\_\allowbreak fine.py}, \texttt{lens\_\allowbreak absorber.py},
\texttt{lens\_\allowbreak flux\_\allowbreak control.py} and, for this revision, \texttt{wt\_\allowbreak r1\_\allowbreak lens.py}
(power balance, flux map, fine sweep, absorber series, as-built), \texttt{wt\_\allowbreak r1\_\allowbreak rays.py},
\texttt{wt\_\allowbreak r1\_\allowbreak cloak.py} (mesh series, wider layer, scattered fields),
\texttt{wt\_\allowbreak r1\_\allowbreak dose\_\allowbreak ms.py} and \texttt{wt\_\allowbreak r1\_\allowbreak panel.py}, all in
\texttt{examples/\allowbreak analysis/}. Table~\ref{tab:verify} lists the lens checks.
\begin{itemize}[itemsep=1pt]
\item \textbf{Mesh.} The lens is re-solved with the optimized field resampled (through
  $(\cos2\theta,\sin2\theta)$, the field being a smooth CS-RBF field) on structured meshes
  with $n_x=130$ (the design mesh: about 9 P1 elements per wavelength of the slowest,
  quasi-shear wave at $\omega=40$, and about 46 per along-fiber qP wavelength), $195$ and
  $260$. At $\omega=40$, $G$ stays within $\WMeshGmin$--$\WMeshGmax$
  and $G_\varphi$ within $\WMeshPmin$--$\WMeshPmax$. The design's sharpest orientation gradient length,
  $L_\theta=0.15$, spans about $5$ elements of the design mesh. The cloak's fixed-design refinement series
  ($h=0.045$ to $0.025$) is discussed in Sec.~\ref{sec:cloak}.
\item \textbf{Damping.} At $\omega=40$ the redirection is independent of the damping
  ($G_\varphi=\WEtaOneP$, $\WEtaTwoP$, $\WEtaFourP$ for $\eta=0.01$, $0.02$, $0.04$), while
  the interior growth falls roughly as $1/\eta$ ($\WEtaOneGr$, $\WEtaTwoGr$, $\WEtaFourGr$: a ratio of $1.7$--$2.0$ per doubling of $\eta$), and so
  does the gain ($\WEtaOneG$, $\WEtaTwoG$, $\WEtaFourG$). This directly confirms that the
  gain beyond $G_\varphi$ is resonant.
\item \textbf{Absorbing layer.} The reflection coefficient is fitted on the centerline of
  the straight-fiber plate, $u_x=Ae^{-ikx}+Be^{ikx}$ with the complex along-fiber qP
  wavenumber, over $0.5<x<3.3$ (fit residual $\le2\%$); $|R|=|B/A|$ at the sponge entrance.
  The paper's sponge (width $0.6$) gives $|R|=\WReflMin$--$\WReflMax$ over
  $\omega\in[35,50]$ (Table~\ref{tab:verify} lists $37.5$--$50$); doubling its width outward, so that the interior is unchanged, gives
  $|R|\le\WReflThickMax$. Table~\ref{tab:verify} lists the lens metrics for both.
\item \textbf{Re-optimization with the better absorber.} The lens is re-optimized from
  straight fibers with the doubled sponge (same objective, MMA settings, 70 iterations and
  mirror-symmetric CS-RBF field; \texttt{lens\_\allowbreak absorber.py}). At $\omega=40$ it
  reaches $G=\WAbsG\times=\WAbsP\times$ ($G_\varphi$) $\times\,\WAbsGr\times$ (growth), with
  $G_\varphi\ge\WAbsMinFortySixty\times$ over $[40,60]$ and a band average of
  $\WAbsBand\times$, equal to that of the original design evaluated with the same absorber
  ($\WAbsBandOld\times$). The original design does not exploit the reflections; the
  reflections raise the redirection metric of any design.
\item \textbf{Panel mesh.} The band design re-evaluated on a finer conforming mesh
  ($h=0.04$ instead of $0.055$; \texttt{panel\_\allowbreak mesh\_\allowbreak check.py}) keeps every
  qualitative result across the band: the sensor at or above $\WPanFineSenMin\times$, the
  payload at or below $\WPanFinePayMax\times$ and the rims at or below $\WPanFineRimMax\times$
  their straight-fiber energies. Sensor and payload ratios change by at most
  $\WPanMeshSenDev\%$; the rim ratios, the smallest energies, change most (by up to
  $\WPanMeshRimDevAbs$ in absolute terms).
\item \textbf{Cloak frequency response on the finer mesh.} The cloak re-optimized on $h=0.03$
  gives $\WCloakFineMinus\times$ and $\WCloakFinePlus\times$ at $\omega=47.5$ and $52.5$
  against $\WCloakFineReopt\times$ at $50$ (\texttt{cloak\_\allowbreak fine\_\allowbreak band.py}),
  as narrowband as the design on the coarser mesh.
\end{itemize}
\begin{table}[!tbp]
\centering\footnotesize
\caption{Numerical verification of the lens (fixed optimized design).}\label{tab:verify}
\IfFileExists{wt_verify_table.tex}{\begin{tabular}{@{}llccc@{}}
\toprule
check & setting & $G$ & $G_\varphi$ & growth\\ \midrule
lens mesh, $\omega=40$ & $n_x=130$ & 78.0 & 5.82 & 13.41\\
lens mesh, $\omega=40$ & $n_x=195$ & 73.2 & 5.72 & 12.80\\
lens mesh, $\omega=40$ & $n_x=260$ & 74.4 & 5.96 & 12.48\\
lens mesh, $\omega=50$ & $n_x=130$ & 5.5 & 6.01 & 0.91\\
lens mesh, $\omega=50$ & $n_x=195$ & 5.4 & 5.95 & 0.90\\
lens mesh, $\omega=50$ & $n_x=260$ & 5.7 & 6.32 & 0.90\\
\midrule
lens damping, $\omega=40$ & $\eta=0.01$ & 135.1 & 5.80 & 23.32\\
lens damping, $\omega=50$ & $\eta=0.01$ & 5.1 & 6.03 & 0.85\\
lens damping, $\omega=40$ & $\eta=0.02$ & 78.0 & 5.82 & 13.41\\
lens damping, $\omega=50$ & $\eta=0.02$ & 5.5 & 6.01 & 0.91\\
lens damping, $\omega=40$ & $\eta=0.04$ & 37.9 & 5.77 & 6.57\\
lens damping, $\omega=50$ & $\eta=0.04$ & 5.9 & 5.94 & 1.00\\
\midrule
lens sponge, $\omega=37.5$ & width 0.6, $|R|=0.31$ & 6.6 & 3.74 & 1.77\\
lens sponge, $\omega=37.5$ & width 1.2, $|R|=0.12$ & 4.0 & 3.00 & 1.32\\
lens sponge, $\omega=40$ & width 0.6, $|R|=0.34$ & 78.0 & 5.82 & 13.41\\
lens sponge, $\omega=40$ & width 1.2, $|R|=0.11$ & 64.1 & 4.38 & 14.62\\
lens sponge, $\omega=45$ & width 0.6, $|R|=0.33$ & 16.3 & 6.54 & 2.49\\
lens sponge, $\omega=45$ & width 1.2, $|R|=0.11$ & 10.3 & 4.56 & 2.26\\
lens sponge, $\omega=50$ & width 0.6, $|R|=0.32$ & 5.5 & 6.01 & 0.91\\
lens sponge, $\omega=50$ & width 1.2, $|R|=0.11$ & 6.7 & 4.61 & 1.46\\
\bottomrule
\end{tabular}
}{\textit{table pending}}
\end{table}

